\begin{abstract}
La información y la desinformación modifican la evolución de una epidemia, y sus escalas temporales distintas vuelven rígido el modelo que las acopla. Se modeló numéricamente el sistema UMZ-SIRS para analizar la estabilidad numérica y la persistencia epidémica bajo separación de escalas. Se compararon Euler explícito, Runge--Kutta de cuarto orden (RK4), Runge--Kutta--Fehlberg 4(5) (RKF45) y Euler hacia atrás con Newton--Raphson y Gauss--Seidel, validando las referencias con tolerancias decrecientes, integradores independientes y el equilibrio analítico; los parámetros se exploraron con barridos y mapas. RK4 mostró orden cuatro, RKF45 requirió cinco veces menos evaluaciones que RK4 para errores de $10^{-9}$ a $10^{-11}$ y los métodos explícitos exigieron pasos menores que 2.04 y 2.82. Gauss--Seidel limitó Euler hacia atrás a pasos de 2.1 a 3.3, límite que desapareció con solución directa. Epidemiológicamente, la desinformación convirtió una epidemia con $R_{00}=0.9$ en endemia ($I=0.088$), el escepticismo extinguió una endemia y la información indujo endemia con $R_{00}<1$ en el 8\,\% y el 41\,\% de los puntos de los mapas. Un integrador adaptativo resolvió la rigidez con menor costo y la información puede invertir el destino de la epidemia.
\end{abstract}
\noindent\textbf{Palabras clave:} UMZ-SIRS; rigidez numérica; Euler hacia atrás; Gauss--Seidel; desinformación.
